% ==================================================================
% id:       func-024
% topic:    פונקציות
% subtopic: פונקציה הפוכה
% solution: full
% ==================================================================
\textbf{נתבונן בפונקציה
$f(x)=\frac{x^2-1}{x}$}
\begin{enumerate}
\item
מה תחום ההגדרה של הפונקציה?
\begin{solbox}
תחום ההגדרה
$D=\{x\neq 0\}$
\end{solbox}
\item
האם הפונקציה מחזורית?
\hint{אם נראה לכם שהפונקציה מחזורית, מצאו מחזור; אחרת חפשו תכונה שלה (אפשר להיעזר בסעיף ד') שסותרת את $f(x+T)=f(x)$ לכל $x$.}
\begin{solbox}
לפי סעיף ד' נראה שהיא עולה ממש כאשר
$x>0$
מכך ניתן להסיק שהיא לא מחזורית.
\end{solbox}
\item
האם הפונקציה זוגית?
\begin{solbox}
הפונקציה לא זוגית כי לדוגמה
$f(2)=\frac{3}{2}, f(-2)=-\frac{3}{2}$
\end{solbox}
\item
באילו תחומים הפונקציה מונוטונית עולה או יורדת?
\hint{כתבו את $f(x)$ כסכום של שני איברים פשוטים, ובדקו את המונוטוניות של כל אחד מהם בכל צד של $0$ בנפרד.}
\begin{solbox}
נשים לב כי
$f(x)=\frac{x^2-1}{x}=x-\frac{1}{x}$
ו-
נשים לב כי עבור
$x>0$
הפונקציה $x$ היא  עולה
ממש
לכן $\frac{1}{x}$
היא יורדת
ממש
לכן $-\frac{1}{x}$
עולה ממש ולכן
$f(x)$
עולה ממש כסכום עולות ממש.
\\
נשים לב כי עבור
$x<0$
הפונקציה $x$ היא  עולה
ממש
לכן $\frac{1}{x}$
היא יורדת
ממש
לכן $-\frac{1}{x}$
עולה ממש ולכן
$f(x)$
עולה ממש כסכום עולות ממש.
\end{solbox}
\item
האם הפונקציה הפיכה? אם לא מצאו תחום בו הפונקציה הפיכה.
\hint{בדקו חד-חד ערכיות על כל תחום ההגדרה, ואם צריך היעזרו בסעיף ד' כדי לבחור תחום מתאים.}
\begin{solbox}
הפונקציה לא חח"ע למשל כי
$f(-1)=f(1)=0$
ולכן לא הפיכה
.
\\
אבל אם נגביל את עצמנו לאחד מהתחומים מימין ל- $0$ או משמאל ל- $
0$
הרי ששם הפונקציה עולה ממש ולכן חח"ע.\\

למשל בתחום
$x>0$
אם נקח את
$B=\{f(x)|x>0 \}$
אז ברור שבתחום הזה $f$
על
$B$
ולכן הפיכה.
\end{solbox}
\end{enumerate}
